\section{No information without disturbance}
\label{sec:no_information_without_disturbance}
%% =========================================================

An \emph{instrument} on $\MN{d}$ is a finite family
$\{T_\alpha:\MN{d}\to\MN{d}\}$ of completely positive maps, the outcome
$\alpha$ occurring with probability $\tr[T_\alpha(\rho)]$ on the input
$\rho$. The instrument leaves the input undisturbed on average when
$\sum_\alpha T_\alpha=\id$. In that case the outcome distribution is the
same for every input, so the measurement returns no information about
$\rho$. This is the proposition ``No information without disturbance'' of
\cite[Chapter~2]{Wolf2012Quantum}.

\begin{lemma}[Domination by the maximally entangled projector]
    \label{lem:le_omega_proj_eq_nonneg_smul}
    \lean{Channel.exists_nonneg_smul_omegaProj_of_le_omegaProj}
    \leanok
    \uses{def:maximally_entangled}
    Let $A$ on $\C^d\otimes\C^d$ be positive semidefinite with
    $A\leq\ketbra{\Omega}{\Omega}$. Then $A=a\ketbra{\Omega}{\Omega}$ for a
    non-negative scalar $a$.
\end{lemma}

\begin{proof}
    \leanok
    \uses{lem:posSemidef_rankone_domination}
    The projector $\ketbra{\Omega}{\Omega}$ is the rank-one matrix built
    from the vector $\ket{\Omega}$, so this is
    Lemma~\ref{lem:posSemidef_rankone_domination} with $c=1$ and
    $\psi=\Omega$, read after transporting along a bijection between the
    pair index set of the bipartite space $\C^d\otimes\C^d$ and a single
    index set of the same cardinality $d^2$.
\end{proof}

\begin{theorem}[No information without disturbance]
    \label{thm:no_information_without_disturbance}
    \lean{Channel.exists_nonneg_smul_id_of_isCPMap_of_sum_eq_id}
    \leanok
    \uses{def:cp_map}
    Let $d\geq1$ and let $\{T_\alpha:\MN{d}\to\MN{d}\}$ be a finite family of
    completely positive maps with $\sum_\alpha T_\alpha=\id$. Then for every
    $\alpha$ there is a constant $c_\alpha\geq0$ with
    $T_\alpha=c_\alpha\id$.
\end{theorem}

\begin{proof}
    \leanok
    \uses{def:choi_matrix, def:maximally_entangled,
        thm:cp_iff_choi_posSemidef, thm:choi_matrix_id,
        lem:choi_matrix_injective, lem:le_omega_proj_eq_nonneg_smul}
    Pass to Choi matrices. The Choi matrix of $\id$ is
    $\ketbra{\Omega}{\Omega}$ and the Choi correspondence is linear, so the
    hypothesis reads $\ketbra{\Omega}{\Omega}=\sum_\alpha\tau_\alpha$,
    where $\tau_\alpha$ is the Choi matrix of $T_\alpha$. Each $\tau_\alpha$
    is positive semidefinite because $T_\alpha$ is completely positive, so
    every summand is dominated by the whole sum:
    $\tau_\alpha\leq\ketbra{\Omega}{\Omega}$. The maximally entangled state
    is pure, hence admits no non-trivial decomposition, and
    Lemma~\ref{lem:le_omega_proj_eq_nonneg_smul} gives
    $\tau_\alpha=c_\alpha\ketbra{\Omega}{\Omega}$ with $c_\alpha\geq0$. The
    right-hand side is the Choi matrix of $c_\alpha\id$, so
    Lemma~\ref{lem:choi_matrix_injective} gives $T_\alpha=c_\alpha\id$.
\end{proof}

\begin{theorem}[Normalization of the outcome weights]
    \label{thm:no_information_weights_sum_one}
    \lean{Channel.exists_nonneg_weights_of_isCPMap_of_sum_eq_id}
    \leanok
    \uses{thm:no_information_without_disturbance}
    Under the hypotheses of
    Theorem~\ref{thm:no_information_without_disturbance} the constants
    satisfy $c_\alpha\geq0$ and $\sum_\alpha c_\alpha=1$.
\end{theorem}

\begin{proof}
    \leanok
    \uses{thm:no_information_without_disturbance}
    Evaluate $\sum_\alpha T_\alpha=\id$ at the identity matrix and take the
    trace: the left-hand side gives $d\sum_\alpha c_\alpha$ and the
    right-hand side gives $d$. Divide by $d\geq1$.
\end{proof}

\begin{theorem}[The outcome probability is independent of the input]
    \label{thm:no_information_probability}
    \lean{Channel.exists_nonneg_forall_trace_map_eq_of_isCPMap_of_sum_eq_id}
    \leanok
    \uses{thm:no_information_without_disturbance}
    Under the hypotheses of
    Theorem~\ref{thm:no_information_without_disturbance} there is, for every
    $\alpha$, a constant $c_\alpha\geq0$ with
    $\tr[T_\alpha(\rho)]=c_\alpha$ for every $\rho$ of unit trace. The
    probability of the outcome $\alpha$ therefore does not depend on the
    input, so no information is gained.
\end{theorem}

\begin{proof}
    \leanok
    \uses{thm:no_information_without_disturbance}
    By Theorem~\ref{thm:no_information_without_disturbance},
    $T_\alpha(\rho)=c_\alpha\rho$, whose trace is $c_\alpha\tr(\rho)=c_\alpha$.
\end{proof}

\begin{corollary}[No information without disturbance for an instrument]
    \label{cor:no_information_instrument}
    \lean{Instrument.exists_nonneg_forall_probability_eq_of_total_eq_id}
    \leanok
    \uses{def:instrument, def:instrument_operations}
    Let $\{\Phi_i\}$ be a quantum instrument on $\MN{d}$ with $d\geq1$ whose
    total channel is the identity. Then for every outcome $i$ the probability
    $p_i(\rho)$ is a non-negative constant, the same for every $\rho$ of unit
    trace.
\end{corollary}

\begin{proof}
    \leanok
    \uses{thm:no_information_probability}
    The component maps of an instrument are completely positive and their sum
    is the identity by hypothesis, so
    Theorem~\ref{thm:no_information_probability} applies.
\end{proof}
